\section{Lowest Landau level and zero modes}\label{sec:prelim}

We fix the normalization of the interaction before discussing locality.
The exterior-algebra formulation is useful because both the pinned
Hamiltonians and the particle-loss evolution change particle number.
Throughout the paper, constants are independent of flux and particle
number unless an explicit dependence is indicated. A statement for large
flux means that its threshold is independent of the particle sector.

\subsection{Coherent orbitals and pair annihilators}

Let $q\ge1$ be an integer. The sphere $S_q$ has radius $\sqrt{q/2}$,
area $2\pi q$, and magnetic field one. The lowest Landau level $U_q$ is
the spin-$q/2$ representation of $SU(2)$. We use homogeneous polynomials
of degree $q$ in spinor coordinates $(z_0,z_1)$, with invariant inner
product normalized so that
\[
 e_j=\binom qj^{1/2}z_0^{q-j}z_1^j,
 \qquad 0\le j\le q,
\]
is an orthonormal basis. This normalization identifies $U_q$ unitarily
with the physical lowest Landau level; it does not rescale any operator.
Relative to a chosen north pole, $j=0$ is the north orbital and $j=q$
the south orbital.

Write $k_x$ for the unit coherent vector at $x\in S_q$, defined up to
a phase by the Riesz vector for evaluation at $x$. Its components in
the axis basis have squared moduli
\begin{equation}\label{eq:binomial-coherent}
 |\langle e_j,k_x\rangle|^2
 =\binom qj (1-m/q)^{q-j}(m/q)^j,
 \qquad m=q\sin^2(\theta/2).
\end{equation}
Schur's Lemma and the trace give the resolution
\begin{equation}\label{eq:coherent-resolution}
 \kappa_q\int_{S_q}|k_x\rangle\langle k_x|\,dA(x)=I_{U_q},
 \qquad \kappa_q=\frac{q+1}{2\pi q}.
\end{equation}
In particular the density $\kappa_q$ is uniformly bounded. For a bounded
measurable real function $\varphi$, the Toeplitz operator is
\begin{equation}\label{eq:toeplitz-resolution}
 T_q(\varphi)=\kappa_q\int_{S_q}
       \varphi(x)|k_x\rangle\langle k_x|\,dA(x).
\end{equation}
Indeed the quadratic form on the right is the integral of $\varphi$
times the product of the two evaluated lowest-level sections, which is
the quadratic form of compression of multiplication by $\varphi$.

The fermionic Fock space and number operator are
\[
 \mathcal F_q=\bigoplus_{n=0}^{q+1}\bigwedge^n U_q,
 \qquad \mathcal N\big|_{\wedge^n U_q}=nI.
\]
Increasing wedges of orthonormal vectors have norm one. For
$a\in\bigwedge^k U_q$, let $B(a)$ be the adjoint of left wedge
multiplication by $a$. Thus $B$ is conjugate-linear,
$B(a\wedge b)=B(b)B(a)$, and $c_j=B(e_j)$ satisfy the canonical
anticommutation relations. For a one-particle vector $f$ we also write
$c(f)=B(f)$. For an operator on $\bigwedge^k U_q$, its
lift to Fock space is the linear extension of
\begin{equation}\label{eq:lift}
 \mathcal L_k(|a\rangle\langle b|)=B(a)^*B(b).
\end{equation}
The lift preserves positivity and rotations. It vanishes on sectors
with fewer than $k$ particles and equals its argument on the $k$-particle
sector.

Let $V_q\subset\bigwedge^2 U_q$ be the spin-$(q-1)$ summand. Its
orthonormal axis basis $v_p$, $0\le p\le2q-2$, can be written as
\begin{equation}\label{eq:pair-rows}
 v_{b-1}=Z_{q,b}^{-1/2}
 \sum_{\substack{0\le i<j\le q\\i+j=b}}
 (j-i)\binom qi^{1/2}\binom qj^{1/2}e_i\wedge e_j,
 \quad
 Z_{q,b}=\frac{b(2q-b)}{2(2q-1)}\binom{2q}{b},
\end{equation}
where $1\le b\le2q-1$. The formula follows by lowering the highest
vector $e_0\wedge e_1$. For its normalization, the sum over ordered
$i,j$ with $i+j=b$ is a hypergeometric second moment: its total mass
is $\binom{2q}{b}$ and the mean of $(j-i)^2$ is
$b(2q-b)/(2q-1)$. Dividing by two gives $Z_{q,b}$.
We set $B_p=B(v_p)$. The interaction is
\begin{equation}\label{eq:fock-H}
 H_q=\mathcal L_2(\Pi_{V_q})=
 \sum_{p=0}^{2q-2}B_p^*B_p,
 \qquad H_{n,q}=H_q\big|_{\wedge^n U_q}.
\end{equation}
On $n$ particles this is exactly $\sum_{i<j}P_{ij}^{(1)}$, with
coefficient one for each unordered pair. One way to check the coefficient
is to identify a unit wedge with its normalized antisymmetric tensor:
contraction in two prescribed tensor positions is $B(a)/\sqrt{\binom n2}$,
and summing the $\binom n2$ pair projections gives
\eqref{eq:fock-H}.

We shall also use the wedge maps
\begin{equation}\label{eq:wedge-maps}
 \begin{aligned}
 W_3 &: V_q\otimes U_q\longrightarrow\bigwedge^3 U_q,
       &W_3(a\otimes u)&=a\wedge u,\\
 W_4 &: V_q\otimes V_q\longrightarrow\bigwedge^4 U_q,
       &W_4(a\otimes b)&=a\wedge b.
 \end{aligned}
\end{equation}
In particular $V_q\otimes V_q$ uses ordered pairs. This convention
fixes the coefficient of the four-body term in the normal-ordered square
used in Appendix~\ref{app:retained-blocks}.

\subsection{The zero spaces and the unperturbed input}

The polynomial and second-quantized descriptions of these kernels are
closely related; see~\cite{CS,MONS} for algebraic constructions of
pseudopotential zero modes. Put $Z_{n,q}=\ker H_{n,q}$ and let $P_{n,q}$ be its orthogonal
projection. When all sectors are being considered, write
$Z_q=\ker H_q$ and $P_q$ for the Fock-space zero projection.

\begin{lemma}[Polynomial description]\label{lem:zero-modes}
For $n\ge1$, a vector lies in $Z_{n,q}$ if and only if its polynomial
has the form
\begin{equation}\label{eq:zero-polynomial}
 \prod_{i<j}(z_{i0}z_{j1}-z_{i1}z_{j0})^3\,S(z_1,\ldots,z_n),
\end{equation}
where $S$ is symmetric and homogeneous of degree
$h=q+3-3n$ in each spinor. Consequently $Z_{n,q}=\{0\}$ when
$h<0$, and, when $h\ge0$,
\[
 \dim Z_{n,q}=\binom{n+h}{n}.
\]
At $q=3(N-1)$ the space $Z_{N,q}$ is the Laughlin line, and every
nonempty zero sector has particle number at most $N$.
The vacuum sector $Z_{0,q}$ has dimension one.
\end{lemma}
\begin{proof}
Positivity of the pair terms says that a zero vector has no pair
relative-angular-momentum-one component. An antisymmetric polynomial
of the same homogeneous degree in each of two spinors vanishes when
the spinors coincide. Separate homogeneity extends that vanishing to
proportional spinors, so the irreducible bracket $[z_1,z_2]$ divides it.
Dividing by
the bracket and then setting $z_1=z_2=z$ gives an equivariant map to
the degree-$(2q-2)$ polynomials in $z$, the spin-$(q-1)$ representation.
This map is nonzero on $e_0\wedge e_1$ and hence is onto. Its kernel
is exactly divisibility by the bracket cubed: the divided polynomial
is symmetric and separately homogeneous, and if it vanishes on the
diagonal it is divisible by
the bracket; the resulting quotient is antisymmetric and therefore
divisible by the bracket once more. By the two-spin decomposition,
the map detects precisely the spin-$(q-1)$ summand. Thus absence of
that component is equivalent to cubic divisibility.

It follows that every pair bracket cubed divides the polynomial.
The distinct pair brackets are nonassociate irreducible polynomials;
unique factorization makes their product divide it. Dividing by the
cubic product changes antisymmetry to symmetry and leaves degree
$q-3(n-1)$ in each variable. This proves both directions of
\eqref{eq:zero-polynomial}. A symmetric polynomial with individual
degree $h$ has the symmetrized monomial basis indexed by
$0\le j_1\le\cdots\le j_n\le h$, giving the binomial dimension.
The remaining claims follow at once, with the vacuum treated separately.
\end{proof}

The quantitative unperturbed result used in this paper is the following
Fock-space gap theorem. We recall it with the normalization just fixed;
its proof is in the companion manuscript~\cite[Theorem~1.1 and
Corollary~1.2]{Gap}.

\begin{theorem}[Uniform unperturbed gap]\label{thm:unperturbed}
Let $\gamma_*=4616733319001/10^{14}$. For every
$0<\gamma<\gamma_*$ there is $q_\gamma$ such that
\begin{equation}\label{eq:accepted-fock-gap}
 H_q^2\ge\gamma H_q\qquad(q\ge q_\gamma)
\end{equation}
on all of $\mathcal F_q$. In particular, for all sufficiently large
$N$ and $q=3(N-1)$,
\begin{equation}\label{eq:accepted-filled-gap}
 H_{N,q}\ge\frac1{25}(I-P_{N,q}).
\end{equation}
\end{theorem}

The all-sector statement gives separation of the spectrum from zero.
We use it for zero-space transport, including particle numbers below
Laughlin filling. The extra four-body comparison needed for energy
descent is proved separately in Appendix~\ref{app:retained-blocks}; it
is stronger than \eqref{eq:accepted-fock-gap}.

\subsection{Geometric scales}

An axis on $S_q$ gives latitude coordinate $m=q\sin^2(\theta/2)$
and longitude $\phi$. In these coordinates $dA=dm\,d\phi$.
We use
\begin{equation}\label{eq:geometric-scales}
 u=1+m,\qquad v=1+q-m,\qquad
 w=\sqrt{\frac{uv}{q+1}}.
\end{equation}
Here $v$ is an end-distance weight; the potential is always denoted
by $\varphi$. The width $w$ converts index distance to physical distance,
including in the polar caps. Section~\ref{sec:locality} proves the
geometric estimates that make these scales useful for local operators.
